\titleframe

% Contoh Kode Slide
%\section{Introduction}
%\begin{frame}{\insertsectionhead}
%	\framesubtitle{A short introduction to Trigon}
%	\themename is a modern, elegant and versatile theme for Beamer, inspired by
%	the
%	\href{https://github.com/matze/mtheme}{\textsc{metropolis} theme} from Matthias
%	Vogelgesang.
%	\vfill
%	\themename comes with lots of nice extra features
%	\begin{itemize}
%		\item Multiple style variations for title, section and normal slides
%		\item Simple customization of theme colors
%		\item Lots of convenient options to tweak the design
%	\end{itemize}
%\end{frame}

% Contoh Kode Tabel
%  \begin{table}[H]
% \centering
% \caption{A nice table example}
% \begin{tabular}{@{} lccc @{}}
%    \toprule
%    & \textbf{Velocity} & \textbf{Angle}  & \textbf{Vertical force} \\
%    & $U$ & $\alpha$  & $F_z$ \\
%    & [m/s] & [$^\circ$]  & [N] \\
%    \midrule
%    2D simulation  & 9 & 2 & 9.23 \\
%    3D simulation  & 10.0 & 3 & 15.039 \\
%    Experiment A   & 11.31 & 2.5 & 13.2 \\
%    Experiment B   & 11.26 & 2.7 & 12.6 \\
%    Experiment C   & 11.33 & 2.47 & 13.6 \\
%    \bottomrule
%  \end{tabular}
%\end{table}

% Contoh Kode Gambar
%\begin{figure}[ht!]
%	\begin{subfigure}[b]{0.3\textwidth}
%		\frame{\includegraphics[width=\textwidth]{layout_example-03.jpg}}
%		\caption*{plain}
%	\end{subfigure}
%	\hspace{\fill}
%	\begin{subfigure}[b]{0.3\textwidth}
%		\frame{\includegraphics[width=\textwidth]{layout_example-02.jpg}}
%		\caption*{style1}
%	\end{subfigure}
%	\hspace{\fill}
%	\begin{subfigure}[b]{0.3\textwidth}
%		\frame{\includegraphics[width=\textwidth]{layout_example-01.jpg}}
%		\caption*{style2 (default)}
%	\end{subfigure}
%\end{figure}

% Contoh Kode Blocks
% \begin{block}{Regular block}
%	Just a regular block
%\end{block}
%\begin{alertblock}{Alert block}
%	Some important thing
%\end{alertblock}
%\begin{exampleblock}{Example block}
%	No difference with regular block to avoid excessive distraction
%\end{exampleblock}

\setbeamertemplate{frame footer}{Alauddin Maulana Hirzan, M. Kom - Universitas Semarang}

\section{Otomatisasi \textit{Internet of Things}}
\begin{frame}{\insertsectionhead}
	\framesubtitle{Mengenal Otomatisasi}
	\justifying
	Menurut \textbf{ibm.com}:
	\vfill
	\begin{quote}
		Penerapan teknologi, program, robotika, atau proses untuk mencapai hasil dengan masukan manusia yang minimal
	\end{quote}
	\vfill
	Otomasi adalah istilah untuk aplikasi teknologi di mana input manusia diminimalkan. Ini termasuk otomatisasi proses bisnis (BPA), otomatisasi TI, aplikasi pribadi seperti otomatisasi rumah dan banyak lagi.
\end{frame}

\begin{frame}{\insertsectionhead}
	\framesubtitle{\textit{Mass Rapid Transit?}}
	\justifying
	\begin{figure}[ht!]
		\begin{subfigure}[b]{0.5\textwidth}
			\frame{\includegraphics[width=\textwidth]{mrt.jpg}}
		\end{subfigure}
	\end{figure}
\end{frame}

\begin{frame}{\insertsectionhead}
	\framesubtitle{\textit{Mass Rapid Transit?}}
	\justifying
	Apa yang hebat dari MRT yang sepenuhnya otomatis dan kereta tanpa pengemudi?
	\begin{itemize}
		\item Keamanan yang ditingkatkan: Setiap kereta dipantau dari jarak jauh dari Pusat Kontrol Operasi dan Ruang Kontrol Stasiun di mana staf secara aktif mencari kerusakan
		\item Peningkatan keandalan dan efisiensi : Untuk mencegah gangguan layanan apa pun, Sistem Kontrol dan Pengawasan Terpadu aktif mencari pemecahan masalah untuk setiap kelainan operasional dan memberi tahu tim pemeliharaan secara real time
		\item Ramah perawatan : Semua peralatan dirancang dalam bentuk modular, memungkinkan pekerjaan penggantian atau perbaikan dilakukan dengan relatif mudah dan gangguan minimal.
	\end{itemize}
\end{frame}

\begin{frame}{\insertsectionhead}
	\framesubtitle{Prinsip Dasar Otomatisasi}
	\justifying
	Otomatisasi dasar mengambil tugas-tugas sederhana yang belum sempurna dan mengotomatiskannya. Tingkat otomatisasi ini adalah tentang mendigitalkan pekerjaan dengan menggunakan alat untuk merampingkan dan memusatkan tugas-tugas rutin, seperti menggunakan sistem pesan bersama
	\vfill
	\begin{block}{Info}
		Intruksi-instruksi sederhana dapat diambil oleh otomatisasi. Sehingga dapat meningkatkan efisiensi dan efektivitas.
	\end{block}
\end{frame}

\begin{frame}{\insertsectionhead}
	\framesubtitle{Tipe Otomatisasi Secara Umum}
	\justifying
	\begin{itemize}
		\item \textbf{Otomatisasi Dasar} : Memusatkan kegiatan rutin secara otomatis seperti berkirim pesan menggunakan Bot/Alat lainnya
		\item \textbf{Otomatisasi Proses} :  Mengatur proses yang biasanya dilakukan oleh software maupun aplikasi tersendiri
		\item \textbf{Otomatisasi AI} : Menggunakan AI untuk melakukan otomatisasi
	\end{itemize}
\end{frame}

\begin{frame}{\insertsectionhead}
	\framesubtitle{Otomatisasi Dasar \#1}
	\justifying
	Biasanya digunakan untuk pekerjaan-pekerjaan sederhana seperti pemasangan baut mobil yang ada di \textit{assembly line}.
	\vfill
	\begin{figure}[ht!]
		\begin{subfigure}[b]{0.53\textwidth}
			\frame{\includegraphics[width=\textwidth]{assembly-line.jpg}}
		\end{subfigure}
	\end{figure}
\end{frame}

\begin{frame}{\insertsectionhead}
	\framesubtitle{Otomatisasi Dasar \#2}
	\justifying
	Teknologi ini sangat membantu produksi yang lebih cepat dan efisien, selain itu juga lebih \textit{cost effective}  dibandingkan menggunakan pekerja biasa.
	\vfill
	Bisnis masih mempertimbangkan atau mempelajari apakah mereka akan beradaptasi dengan otomatisasi. Karena tidak semua bisnis perlu diotomatisasi. Produksi otomatis ini dapat sangat diuntungkan oleh perusahaan besar karena mereka membutuhkan produksi yang cepat dan akurat.
\end{frame}

\begin{frame}{\insertsectionhead}
	\framesubtitle{Otomatisasi Proses \#1}
	\justifying
	Sebagai contoh ada Pencucian Mobil. Yang di aman semua prosesnya sepenuhnya otomatis:
	\vfill
	\begin{itemize}
		\item Pelanggan memilih jenis cucian yang diinginkan
		\item Sistem meminta pembayaran dari pelanggan
		\item Mengidentifikasi saat mobil berada di posisi yang tepat menggunakan sensor
		\item Menggunakan berbagai sensor untuk mengidentifikasi tinggi dan ukuran mobil
		\item Menjalankan program yang dipilih sebelumnya: membilas, mencuci dengan sabun dan sikat, aplikasi lilin dan mengeringkan mobil.
		\item Setelah selesai, menyarankan pelanggan untuk keluar dari tempat cuci mobil.
	\end{itemize}
\end{frame}

\begin{frame}{\insertsectionhead}
	\framesubtitle{Otomatisasi Proses \#2}
	\justifying
	Contoh lain bisa berupa:
	\begin{itemize}
		\item Pembaruan software sistem operasi
		\item Pendaftaran pasien di rumah sakit termasuk verifikasi ID
		\item Online Shopping
		\item Pembayaran Gaji
		\item Pengembangan Software
		\item Monitoring Server dan Klien
	\end{itemize}
\end{frame}

\begin{frame}{\insertsectionhead}
	\framesubtitle{Otomatisasi AI \#1}
	\justifying
	Sesuai dengan namanya, otomatisasi ini dilakukan oleh \textit{Artificial Intelligence}. Sehingga memerlukan teknologi yang tinggi dibandingkan otomatisasi dasar. 
	\vfill
	\begin{block}{Info}
		\textit{Self-Driving Car}  merupakan contoh myata dari otomatisasi menggunakan AI, teknologi ini sangat memerlukan komputasi AI untuk melihat adanya orang yang menyeberang jalan dan menentukan aksi secara independen
	\end{block}
\end{frame}

\begin{frame}{\insertsectionhead}
	\framesubtitle{Otomatisasi AI \#2}
	\justifying
	\begin{figure}[ht!]
		\begin{subfigure}[b]{0.7\textwidth}
			\frame{\includegraphics[width=\textwidth]{tesla.jpeg}}
		\end{subfigure}
	\end{figure}
\end{frame}

\begin{frame}{\insertsectionhead}
	\framesubtitle{Otomatisasi AI \#3}
	\justifying
	\begin{figure}[ht!]
		\begin{subfigure}[b]{0.7\textwidth}
			\frame{\includegraphics[width=\textwidth]{tesla-vision.png}}
		\end{subfigure}
	\end{figure}
\end{frame}

\begin{frame}{\insertsectionhead}
	\framesubtitle{Otomatisasi dan \textit{Internet of Things} \#1}
	\justifying
	Dengan teknologi IoT yang sudah semakin canggih, ditambah dengan adanya otomatisasi membuat kemungkinan untuk banyak terjadi. Otomatisasi menghubungkan komponen komponen yang ada dan memastikan semua proses berjalan dengan baik tanpa gangguan manusia.
	\vfill
	\begin{block}{Info}
		Lemahnya komputasi IoT menjadikan perangkat IoT tidak cocok untuk otomatisasi berbasiskan AI. Namun perangkat IoT cocok untuk otomatisasi dasar tanpa komputasi berat.
	\end{block}
\end{frame}

\begin{frame}{\insertsectionhead}
	\framesubtitle{Otomatisasi dan \textit{Internet of Things} \#2}
	\justifying
	Komponen-komponen yang diperlukan agar perangkat \textit{Internet of Things} bisa digunakan untuk otomatisasi:
	\vfill
	\begin{itemize}
		\item \textbf{Sensor} untuk mengumpulkan data dari lingkungan agar sistem IoT dapat memprosesnya
		\item \textbf{Koneksi} dan identifikasi dari perangkat ke sistem IoT melalui alamat IP
		\item \textbf{Aktuator} yang memungkinkan perangkat untuk mengambil tindakan berdasarkan data dari sensor mereka sendiri dan umpan balik jaringan
	\end{itemize}
\end{frame}

\begin{frame}{\insertsectionhead}
	\framesubtitle{Otomatisasi dan \textit{Internet of Things} \#3}
	\justifying
	\begin{itemize}
		\item \textbf{Gateway IoT} untuk menjembatani data dari berbagai perangkat untuk mencapai cloud. Ini juga menerjemahkan protokol perangkat menjadi satu protokol standar dan menyaring data yang tidak perlu dari perangkat.
		\item \textbf{Cloud}, tempat semua data dari perangkat IoT dikumpulkan dan diproses oleh perangkat lunak
		\item \textbf{Antarmuka pengguna}, dari mana pengguna mendapatkan data dari perangkat sehingga mereka dapat membuat perintah yang diperlukan yang harus dijalankan perangkat
	\end{itemize}
\end{frame}

\begin{frame}{\insertsectionhead}
	\framesubtitle{Otomatisasi dan \textit{Internet of Things} \#4}
	\justifying
	\begin{figure}[ht!]
		\begin{subfigure}[b]{0.53\textwidth}
			\frame{\includegraphics[width=\textwidth]{iot-home.png}}
		\end{subfigure}
	\end{figure}
\end{frame}

\begin{frame}{\insertsectionhead}
	\framesubtitle{Otomatisasi dan \textit{Internet of Things} \#5}
	\justifying
	Manfaat IoT yang dapat dirasakan:
	\vfill
	\begin{itemize}
		\item \textbf{Peningkatan waktu kerja} : Uptime mengacu pada berapa lama bisnis buka atau berjalan.
		\item \textbf{Efisiensi operasional} : Dalam pasokan dan manajemen aset, otomatisasi dan IoT menghadirkan kecerdasan melalui sensor dan teknologi pemindaian.
		\item \textbf{Tingkatkan Keselamatan} : Aset di tempat kerja jika dilengkapi dengan sensor, dapat mendeteksi kegagalan suku cadang yang akan datang dan memperingatkan personel pemeliharaan untuk menghindari cedera karyawan dari keadaan yang dapat dicegah.
	\end{itemize}
\end{frame}

\begin{frame}{\insertsectionhead}
	\framesubtitle{Otomatisasi dan \textit{Internet of Things} \#6}
	\justifying
	\begin{itemize}
		\item \textbf{Manfaat kesehatan} : Perangkat IoT dapat membantu dokter dan tenaga medis lainnya memantau pasien mereka dari jarak jauh, terutama jika mereka dalam kondisi berbahaya. 
		\item \textbf{Manfaat di bidang pertanian} : Petani — terutama yang memiliki operasi skala besar — dapat menerapkan solusi pertanian cerdas untuk memantau ladang tanaman menggunakan sensor dan mengotomatiskan sistem irigasi dari mana saja. 
		\item \textbf{Manfaat bagi konsumen} : Konsumen juga memiliki banyak keuntungan dari otomatisasi dan IoT. Dari perangkat yang dapat dikenakan, seperti jam tangan pintar atau pelacak kebugaran, yang membantu pengguna lebih memahami kesehatan mereka
	\end{itemize}
\end{frame}

\begin{frame}{\insertsectionhead}
	\framesubtitle{Masalah Otomatisasi \textit{Internet of Things} \#1}
	\justifying
	\begin{itemize}
		\item 1. Masalah Keamanan
		\item Seperti semua perangkat komputasi, keamanan akan menjadi masalah yang lebih besar karena lebih banyak orang menggunakan perangkat rumah pintar. Tentu saja, akan ada berbagai masalah keamanan yang muncul dan menjamurnya perangkat lunak dan perangkat keamanan rumah pintar.
	\end{itemize}
	\vfill
	\begin{block}{Info}
		Perangkat dapat diamankan secara fisik maupun digital. Menggunakan kemasan anti maling maupun sistem IDS
	\end{block}
\end{frame}

\begin{frame}{\insertsectionhead}
	\framesubtitle{Masalah Otomatisasi \textit{Internet of Things} \#2}
	\justifying
	\begin{itemize}
		\item 2. Biaya: Sangat mahal
		\item Meskipun banyak perangkat rumah pintar sekarang terjangkau bagi banyak orang, tetapi tetap saja sangat mahal untuk melengkapi rumah dengan perangkat pintar.
	\end{itemize}
	\vfill
	\begin{block}{Info}
		Biaya dapat ditekan dengan membeli \textit{board} alternatif yang memiliki kemampuan yang sama namun dengan harga yang lebih murah
	\end{block}
\end{frame}

\begin{frame}{\insertsectionhead}
	\framesubtitle{Masalah Otomatisasi \textit{Internet of Things} \#3}
	\justifying
	\begin{itemize}
		\item 3. Penerimaan masyarakat
		\item Banyak teknologi baru terkadang dianggap tidak perlu oleh beberapa orang
	\end{itemize}
	\vfill
	\begin{block}{Info}
		Pendekatan yang berbeda dapat digunakan untuk memperkenalkan teknologi ke masyarakat sekitar.
	\end{block}
\end{frame}

\setbeamertemplate{background} 
{
	\includegraphics[width=\paperwidth,height=\paperheight]{thank-you.jpg}
}
\begin{frame}[plain]
\end{frame}